\begin{helper}[Disjoint component reduces to the repeated-suffix case]
Let \((\varphi,T)\) be a minimum bad pair with
\(2<q:=\rho(T)<n:=|E(T)|\).  Let \(\beta\) be missing at
\(y_i,y_n\), let \(\varepsilon\) be missing at \(y_m\), where
\(i,m\le q-2\), and let \(\gamma\) be missing at \(y_q\).  Suppose
\(\gamma\ne\beta\), \(\varepsilon\ne\gamma\), and
\(\varepsilon\ne\beta\), and suppose \(\gamma,\varepsilon\) are both
unused through \(y_q\).  If the \((\varepsilon,\beta)\)-component
rooted at \(y_n\) avoids every displayed vertex of \(T_{q-1}\), then
this configuration is impossible.
\end{helper}
\begin{proof}
Unfold \noderef{TashTreeMinimumBadPair} and take its bad-pair
conjunct.  By \noderef{TashTreeBadPair} it contains a \(k\)-triple, and
\noderef{KTriple} says that \(\varphi\) is proper.  This is the
properness used below for the swap-missing theorem.

Write \(Q\) for the component and
\(\varphi_1\) for the exact swap from
\noderef{TwoColorComponentSwap}.  We first verify the edge-node
hypotheses needed to apply
\noderef{TashTreeDisjointComponentNormalization} at the cutoff
\(p=q-1\).

Fix a formal edge index \(r<q-1\).  By the attachment field of
\noderef{TreeSequence}, the exact edge node for \(T.edge(r)\) has a
displayed endpoint whose index is at most \(r+1\le q-1\), and that
vertex is an entry of the prefix \(T_{q-1}\) by
\noderef{TreeSequencePrefix}, so its vertex-node avoids \(Q\) by
hypothesis.  If the edge has color
\(\varepsilon\) or \(\beta\), then
\noderef{TwoColorIncidenceGraph} joins its edge-node to this endpoint;
membership of the edge-node in \(Q\) would therefore force membership
of the endpoint, a contradiction.  If the edge is uncolored or has a
color outside \(\{\varepsilon,\beta\}\), the same incidence definition
gives its edge-node no neighbors at all.  Such an isolated edge-node
cannot lie in the component \noderef{TwoColorComponent} rooted at the
distinct vertex-node \(y_n\).  Thus every edge node with \(r<q-1\)
avoids \(Q\).  Notice that this argument explicitly covers off-pair
edge nodes rather than assuming they carry an affected color.

The early witnesses satisfy \(m,i\le q-2<p+1\), and all displayed
vertices through \(p=q-1\), which are exactly the entries of that
prefix by \noderef{TreeSequencePrefix}, avoid \(Q\).  We may therefore apply
\noderef{TashTreeDisjointComponentNormalization} with
\(\alpha=\varepsilon\), \(\beta=\beta\), \(a=m\), \(b=i\), and the
cutoff \(p=q-1\).  It returns a minimum bad pair
\((\varphi_1,T)\), preserves every exact edge at formal index below
\(q-1\), and makes \(\varepsilon\) missing at both \(y_m\) and
\(y_n\).

The remaining exact edge needed for the unused-through predicates is
the edge at formal index \(q-1\), namely displayed edge \(e_q\).  Its
attachment data again gives it an earlier endpoint in
\(V(T_{q-1})\); the entry description in
\noderef{TreeSequencePrefix} identifies this as one of the protected
displayed vertices, which avoids \(Q\).  The same affected-color adjacency
argument, and the same isolated off-pair edge-node argument, prove that
this exact edge-node also avoids \(Q\).  Consequently its value is fixed
by the defining cases of \noderef{TwoColorComponentSwap}.  Combined
with the early-edge equalities returned by the normalization, every
formal edge index \(r\) with \(1\le r<q\), namely the exact range
\(e_2,\ldots,e_q\) in \noderef{TreeColorUnusedThrough}, retains a
non-\(\varepsilon\) value.  Thus \(\varepsilon\) remains unused through
\(y_q\).  Since the swap
only exchanges \(\varepsilon\) and \(\beta\), every exact edge having
color \(\gamma\) is fixed.  Applying the same formal-index range from
\noderef{TreeColorUnusedThrough} shows that \(\gamma\) also remains
unused through \(y_q\).

Finally, \(\gamma\) is distinct from both swapped colors.  The
unaffected-color clause of
\noderef{TwoColorComponentSwapMissingBehavior}, applied with the
properness obtained from \noderef{KTriple}, therefore preserves its
missingness at \(y_q\), whether or not that vertex lies in \(Q\).  Under
the new minimum bad pair we now have \(\varepsilon\) missing at the
early vertex \(y_m\) and at \(y_n\), \(\gamma\) missing at \(y_q\),
and both colors unused through \(y_q\).  Since \(m<q<n\) and
\(q=\rho(T)\), these are exactly the hypotheses of
\noderef{TashTreeControlledRepeatedSuffixContradiction}, with
\(\alpha=\varepsilon\), \(\beta=\gamma\), \(i=m\), and \(j=q\).
That result supplies the desired contradiction.  No preservation of the
full no-repetition invariant is needed in this terminal branch.
\end{proof}
